\documentclass[11pt]{article}

\usepackage{amsmath,amssymb}
\usepackage{xurl}
\usepackage[hidelinks]{hyperref}
\setlength{\emergencystretch}{2em}

% Shared commands for notes in docs/paper-gaps/.

\DeclareUnicodeCharacter{2080}{\ifmmode{}_{0}\else\textsubscript{0}\fi}
\DeclareUnicodeCharacter{2081}{\ifmmode{}_{1}\else\textsubscript{1}\fi}
\DeclareUnicodeCharacter{2082}{\ifmmode{}_{2}\else\textsubscript{2}\fi}
\DeclareUnicodeCharacter{2099}{\ifmmode{}_{n}\else\textsubscript{n}\fi}

\newcommand{\Lean}{\textsc{Lean}}
\ExplSyntaxOn
\NewDocumentCommand{\leanid}{m}{
  \ifmmode
    \textsf{\detokenize{#1}}
  \else
    \group_begin:
    \urlstyle{sf}
    \tl_set:Nn \l_tmpa_tl {#1}
    \tl_replace_all:Nnn \l_tmpa_tl {\_} {_}
    \exp_args:NV \path \l_tmpa_tl
    \group_end:
  \fi
}
\ExplSyntaxOff
\providecommand{\tr}{\operatorname{tr}}
\newcommand{\tnleanIssueBase}{https://github.com/LionSR/TNLean/issues/}
\newcommand{\tnleanPrBase}{https://github.com/LionSR/TNLean/pull/}
\newcommand{\tnleanDocBase}{https://sirui-lu.com/TNLean/docs/}
\newcommand{\ghissue}[1]{\href{\tnleanIssueBase#1}{GitHub issue \##1}}
\newcommand{\ghpr}[1]{\href{\tnleanPrBase#1}{PR \##1}}
\newcommand{\doclink}[1]{\href{\tnleanDocBase#1.html}{\path{#1}}}

% Dirac notation and the identity operator, as in blueprint/src/macros/common.tex.
% \providecommand keeps note-local definitions authoritative.
\usepackage{dsfont}
\providecommand{\ket}[1]{|#1\rangle}
\providecommand{\bra}[1]{\langle #1|}
\providecommand{\braket}[2]{\langle #1 | #2 \rangle}
\providecommand{\ketbra}[2]{|#1\rangle\!\langle #2|}
\providecommand{\Id}{\mathds{1}}
\providecommand{\one}{\mathds{1}}
\providecommand{\C}{\mathbb{C}}
\providecommand{\MN}[1]{M_{#1}(\C)}
\providecommand{\MD}{M_D(\C)}

% External referencing for paper sources.  The build helper writes sanitized
% auxiliary files under build/paper-aux/ and excludes duplicated source labels.
\usepackage{xr}

% Import a generated paper auxiliary file with a caller-supplied label prefix.
% The candidate paths support builds from docs/paper-gaps/, the repository root,
% and build/paper-aux/ itself.
\newcommand{\importpaperaux}[2]{%
  \IfFileExists{../../build/paper-aux/#2.aux}{%
    \externaldocument[#1]{../../build/paper-aux/#2}%
  }{%
    \IfFileExists{build/paper-aux/#2.aux}{%
      \externaldocument[#1]{build/paper-aux/#2}%
    }{%
      \IfFileExists{#2.aux}{%
        \externaldocument[#1]{#2}%
      }{}%
    }%
  }%
}

% General external reference macros
\newcommand{\paperref}[2]{\ref{#1-#2}}
\newcommand{\papereqref}[2]{(\ref{#1-#2})}

% CPSV16 source (1606.00608 / MPDO-22-12-17-2.tex) external document and shortcuts
\importpaperaux{cpsv16-}{1606.00608/MPDO-22-12-17-2-tnlean}
\newcommand{\cpsvref}[1]{\paperref{cpsv16}{#1}}
\newcommand{\cpsveqref}[1]{\papereqref{cpsv16}{#1}}

% Verdict marker: \gapnote{<kind>}{<status>}. Kind is the policy
% classification (clarification, local-correction, scope-restriction,
% unfaithful, false-source, open-gap); status is open, wip (elimination
% underway), resolved, or historical. Machine-read by the site index;
% prints a verdict line.
\newcommand{\gapnote}[2]{\par\noindent\textbf{Verdict:} #1 (#2).\par}


\title{The Product-Marginal Reference in the Relative-Entropy Form of
Strong Subadditivity}
\author{}
\date{2026-07-18}

\begin{document}
\maketitle
\gapnote{clarification}{resolved}

\section{The source identity}

For a bipartite state $\omega_{XY}$, let $\omega_X$ and $\omega_Y$ denote
its marginals. Hayden, Jozsa, Petz, and Winter give the identity
\begin{align}
    D(\omega_{XY} \| \omega_X\otimes\omega_Y)
        &=S(\omega_X)+S(\omega_Y)-S(\omega_{XY}).
    \label{eq:hjpw04_ssa_product_marginal_reference_bipartite_identity}
\end{align}
This is Equation~(4) of Ref.~\cite{HJPW2004}. Its specialization to the
bipartition $A\,|\,BC$ is Equation~(5) of the same source. Subtracting the
corresponding identity for $A\,|\,B$ yields Equation~(6):
\begin{align}
    &D(\rho_{ABC} \| \rho_A\otimes\rho_{BC})
      -D(\rho_{AB} \| \rho_A\otimes\rho_B)
    \notag\\
    &\qquad
      =S(\rho_{AB})+S(\rho_{BC})
       -S(\rho_{ABC})-S(\rho_B).
    \label{eq:hjpw04_ssa_product_marginal_reference_product_difference}
\end{align}
Define the channel
\begin{align}
    T
        &=\operatorname{tr}_C.
    \label{eq:hjpw04_ssa_product_marginal_reference_partial_trace_channel}
\end{align}
The second pair of states in
(\ref{eq:hjpw04_ssa_product_marginal_reference_product_difference}) is the
image of the first pair under $T$. Equality in strong subadditivity is
therefore equivalent to equality in relative-entropy data processing for this
pair. The paragraph following Equation~(7) of Ref.~\cite{HJPW2004}
identifies this equivalence. Equations~(4)--(7) appear on page~3 of
arXiv:quant-ph/0304007v2.

The statement requires no strict-positivity hypothesis. For singular
marginals, relative entropy is interpreted on the relevant support. The
support of a bipartite state is contained in the tensor product of the
supports of its marginals, so both relative entropies in
(\ref{eq:hjpw04_ssa_product_marginal_reference_product_difference}) are
finite.

\section{Equivalent maximally mixed criterion}

Let $d_A$ be the dimension of system $A$. The current relative-entropy theory
also proves
\begin{align}
    &D\!\left(\rho_{ABC} \middle\|
      \frac{\Id_A}{d_A}\otimes\rho_{BC}\right)
      -D\!\left(\rho_{AB} \middle\|
      \frac{\Id_A}{d_A}\otimes\rho_B\right)
    \notag\\
    &\qquad
      =S(\rho_{AB})+S(\rho_{BC})
       -S(\rho_{ABC})-S(\rho_B).
    \label{eq:hjpw04_ssa_product_marginal_reference_mixed_difference}
\end{align}
The two relative entropies evaluate to
\begin{align}
    D\!\left(\rho_{ABC} \middle\|
      \frac{\Id_A}{d_A}\otimes\rho_{BC}\right)
        &=\log d_A+S(\rho_{BC})-S(\rho_{ABC}),
    \label{eq:hjpw04_ssa_product_marginal_reference_mixed_first_evaluation}\\
    D\!\left(\rho_{AB} \middle\|
      \frac{\Id_A}{d_A}\otimes\rho_B\right)
        &=\log d_A+S(\rho_B)-S(\rho_{AB}).
    \label{eq:hjpw04_ssa_product_marginal_reference_mixed_second_evaluation}
\end{align}
The terms $\log d_A$ cancel. Hence
(\ref{eq:hjpw04_ssa_product_marginal_reference_mixed_difference}) gives the
same equality criterion as
(\ref{eq:hjpw04_ssa_product_marginal_reference_product_difference}) without
adding a state hypothesis.

The maximally mixed-reference criterion is formalized by
\leanid{isSSAEquality_iff_maximally_mixed_reference_data_processing_eq} in
\path{TNLean/Channel/Schwarz/SSAEqualityDPI.lean}. The declaration
\leanid{SSAPosDef.traceC_ABC_kronecker_traceA_ABC} in
\path{TNLean/Channel/Schwarz/StrongSubadditivityPosDef.lean} proves that the
reference on $A\otimes B$ is the partial-trace image of the reference on
$A\otimes B\otimes C$.

This maximally mixed theorem is a hypothesis-free equivalent formulation with
a different reference state. The exact product-marginal formulation of
Ref.~\cite{HJPW2004} is also proved.

\section{Singular product-marginal evaluation}

The formal relative entropy is the total trace-log expression
\begin{align}
    D(\rho\,\| \sigma)
        &=\operatorname{Re}\tr
          (\rho(\log\rho-\log\sigma)).
    \label{eq:hjpw04_ssa_product_marginal_reference_total_trace_log}
\end{align}
For positive definite tensor factors, the logarithm of a tensor product
splits in the usual way. For positive semidefinite factors $A$ and $B$, let
$P_A$ and $P_B$ denote their support projections. The singular tensor
logarithm satisfies
\begin{align}
    \log(A\otimes B)
        &=(\log A)\otimes P_B
          +P_A\otimes(\log B).
    \label{eq:hjpw04_ssa_product_marginal_reference_singular_tensor_log}
\end{align}
Here the logarithm is extended by zero on the kernel, consistently with
(\ref{eq:hjpw04_ssa_product_marginal_reference_total_trace_log}).

Simultaneous diagonalization of the tensor factors reduces
(\ref{eq:hjpw04_ssa_product_marginal_reference_singular_tensor_log}) to the
scalar identity
\begin{align}
    \log(ab)
        &=\log(a)\mathbf 1_{b\ne0}
          +\mathbf 1_{a\ne0}\log(b)
    \label{eq:hjpw04_ssa_product_marginal_reference_scalar_log_identity}
\end{align}
for nonnegative eigenvalues $a$ and $b$, with each logarithm extended by zero
at the origin. If $a,b>0$, this is the ordinary logarithm identity. If either
eigenvalue vanishes, both sides vanish. This spectral calculation justifies
the singular product-marginal evaluation used in Equation~(4) of
Ref.~\cite{HJPW2004}; that source does not state the tensor-logarithm formula
verbatim. The declaration
\leanid{Matrix.log_kronecker_posSemidef} formalizes
(\ref{eq:hjpw04_ssa_product_marginal_reference_singular_tensor_log}).

For every positive semidefinite bipartite operator $\omega_{XY}$, the
resulting identity is
\begin{align}
    D(\omega_{XY} \| \omega_X\otimes\omega_Y)
        &=S(\omega_X)+S(\omega_Y)-S(\omega_{XY}).
    \label{eq:hjpw04_ssa_product_marginal_reference_product_marginal_evaluation}
\end{align}
The lifted support projections
$P_X\otimes\Id_Y$ and $\Id_X\otimes P_Y$ both fix $\omega_{XY}$. Pairing
(\ref{eq:hjpw04_ssa_product_marginal_reference_singular_tensor_log}) with
$\omega_{XY}$ removes the support projections, and the two partial-trace
adjoint identities yield
(\ref{eq:hjpw04_ssa_product_marginal_reference_product_marginal_evaluation}).
The declaration
\leanid{quantumRelativeEntropy_product_marginals} formalizes this result,
which is precisely Equation~(4) of Ref.~\cite{HJPW2004}.

Applying
(\ref{eq:hjpw04_ssa_product_marginal_reference_product_marginal_evaluation})
to the bipartitions $A\,|\,BC$ and $A\,|\,B$ proves Equations~(5)--(6) of
Ref.~\cite{HJPW2004}, and therefore proves
(\ref{eq:hjpw04_ssa_product_marginal_reference_product_difference}).
Rearrangement gives the exact biconditional between equality in strong
subadditivity and equality in product-reference data processing. The
declaration
\leanid{isSSAEquality_iff_product_marginal_data_processing_eq} formalizes
this biconditional. This step uses neither a recovery map nor the Petz equality
theorem.

\section{Downstream recovery corollary}

The preceding equality criterion feeds the singular-support Petz theorem.
The product-marginal support inclusion is
\begin{align}
    \ker(\rho_A\otimes\rho_{BC})
        &\subseteq\ker\rho_{ABC}.
    \label{eq:hjpw04_ssa_product_marginal_reference_recovery_support}
\end{align}
It is formalized by \leanid{Matrix.product_marginal_support}. Consequently,
equality in strong subadditivity implies
\begin{align}
    \mathcal R^{\mathrm{raw}}_{\rho_A\otimes\rho_{BC}}(\rho_{AB})
        &=\rho_{ABC},
    \label{eq:hjpw04_ssa_product_marginal_reference_raw_recovery}\\
    (\operatorname{id}_A\otimes
      \widehat{\mathcal R}_{\rho_{BC}})(\rho_{AB})
        &=\rho_{ABC}.
    \label{eq:hjpw04_ssa_product_marginal_reference_completed_local_recovery}
\end{align}
The declarations
\leanid{Matrix.partialTraceRightPetzMap_product_marginal_recovery_of_isSSAEquality}
and
\leanid{Matrix.idTensor_partialTraceRightPetzChannel_traceC_ABC_eq_of_isSSAEquality}
formalize Theorem~3 and Equations~(8), (10), and (11) of
Ref.~\cite{HJPW2004} for the exact product reference. The identity-tensored
statement is restricted to the supported input. It does not assert a global
factorization of the generic completed singular product-reference channel.

\section{Verdict}

The clarification is resolved. On the full density-matrix domain, equality in
strong subadditivity is equivalent to saturation of partial-trace data
processing for the exact product-marginal pair used by
Hayden--Jozsa--Petz--Winter \cite{HJPW2004}. The maximally mixed-reference
theorem remains a separate equivalent criterion.

\begin{thebibliography}{9}

\bibitem{HJPW2004}
  P.~Hayden, R.~Jozsa, D.~Petz, and A.~Winter,
  \emph{Structure of states which satisfy strong subadditivity of quantum
  entropy with equality},
  Communications in Mathematical Physics \textbf{246} (2004), 359--374;
  arXiv:quant-ph/0304007v2.

\end{thebibliography}

\end{document}
